%% Moved verbatim from the main text for the Digital Discovery version (2026-10-25).
%% Labels of moved floats carry an s2: prefix so they do not clash with the main text.

\section{Fixed properties, numerics and the Damk\"ohler number}

\paragraph{Fixed properties} Every run uses \(L = \nDcL\)\,m, \(D_{\mathrm{in}} =
\nDcDin\)\,m, \(\rho_g = \nDcRhog\)\,kg\,m\(^{-3}\), \(C_{pg} = \nDcCp\) and \(C_{ps} =
\nDcCps\)\,J\,kg\(^{-1}\)\,K\(^{-1}\), \(k_z = \nDcKz\)\,W\,m\(^{-1}\)\,K\(^{-1}\) and
\(h_w = \nDcHw\)\,W\,m\(^{-2}\)\,K\(^{-1}\). Axial dispersion follows the packed-bed
correlation \(D_L = 0.7D_m + 0.5\,d_p v/\varepsilon_t\) with \(D_m =
\nDcDmDisp\times10^{-5}\)\,m\(^2\)\,s\(^{-1}\). In the larger dataset the kinetic
coefficient is not sampled but derived from the particle by the Glueckauf relation,
\begin{equation}
k_{\mathrm{LDF}} = \frac{15\,\varepsilon_p D_m/\tau_p}{(d_p/2)^2\,\rho_p\,K},
\qquad K = \frac{q^{*}(c_{\mathrm{in}},T_{\mathrm{in}})}{c_{\mathrm{in}}},
\label{s2:eq:glueckauf}
\end{equation}
with pellet porosity \(\varepsilon_p = \nDcGlEpsp\), tortuosity \(\tau_p = \nDcGlTau\)
and \(D_m = \nDcGlDm\times10^{-5}\)\,m\(^2\)\,s\(^{-1}\). The generator therefore uses
two slightly different values of the molecular diffusivity, one in each correlation;
we report both rather than reconcile them after the data were generated.

\paragraph{Numerics and storage} Equations~\eqref{eq:mass}--\eqref{eq:energy} are
discretised by the method of lines on \nGridNz\ cell-centred finite volumes
(first-order upwind advection, central dispersion and conduction) and integrated with
SciPy's implicit backward-differentiation-formula (BDF) method at relative and absolute tolerances
\(10^{\nDcRtolExp}\) and \(10^{\nDcAtolExp}\) to \nDcSnapshots\ output times;
Section~\ref{sec:model} verifies the solver. Each run is integrated to
\nDcHorizon\ stoichiometric times; if the outlet has not reached the acceptance level
the horizon is extended to \nDcHorizonExt\ stoichiometric times in turn, which
\nDcExtendedPct\,\% of the accepted runs of the larger dataset needed. A draw is
rejected before solving if its feed loads too little or too much of capacity, its
Henry constant is not finite and positive, its isothermal mass-transfer zone is not
resolved by the grid, or its stoichiometric time is implausible; \nDataNrejected\ of
the larger dataset's draws were rejected and \nDataNsamples\ accepted. Fields are
stored as \(c/c_{\mathrm{in}}\), \(q/q_{\max}\) and \(T/T_{\mathrm{in}}\) on a uniform
grid in \(z/L\) and \(t/t_{\mathrm{final}}\), at \(\nDcStoreV\times\nDcStoreV\) for the
larger dataset and \(\nDcStoreLegZ\times\nDcStoreLegT\) for the legacy one.

\paragraph{Damk\"ohler number} The Damk\"ohler number used throughout is computed per
sample as \(\mathrm{Da} = k_{\mathrm{LDF}}\,t_{\mathrm{final}}\), and rung L6 regresses
on its per-material median. Because the horizon is a multiple of the stoichiometric
time that depends on whether it was extended, Da carries that multiple as well as the
kinetic ratio \(k_{\mathrm{LDF}}t_{\mathrm{st}}\); for the \nDcExtendedPct\,\% of runs
that were extended it is larger than the kinetic ratio by the extension factor. We keep the
definition the pre-registration used and state this property rather than redefine
the axis after the result.

\section{The parameter space}

\begin{table}[t]
\centering
\caption{The parameter space of the larger dataset. A ``material'' here is the upper
block: an idealised framework pellet described by its isotherm (capacity, step
position, step steepness, the fraction of capacity on the low-humidity sites, heat)
and by pellet size, density and packing. Capacity \nDcQmaxLo--\nDcQmaxHi\,mol\,kg\(^{-1}\)
is \nDcQmaxGgLo--\nDcQmaxGgHi\,g of water per g. The step is re-anchored at each feed
temperature so that it sits at the stated relative humidity, so one material's
isotherm is not the same curve at every condition. All eleven parameters are sampled by
scrambled Sobol and mapped linearly; every surrogate except rung L6 receives exactly
these values, standardised on its own fitting set, and L6 receives observable
descriptors instead.}
\label{s2:tab:params}
\begin{tabular}{llll}
\toprule
parameter & symbol & units & range \\
\midrule
\multicolumn{4}{l}{\emph{material}} \\
saturation capacity & \(q_{\max}\) & mol\,kg\(^{-1}\) & \nDcQmaxLo--\nDcQmaxHi \\
heat of adsorption & \(\Delta H\) & kJ\,mol\(^{-1}\) & \nDcDHLo\ to \nDcDHHi \\
step position (relative humidity) & \(\mathrm{RH}_{\mathrm{step}}\) & -- & \nDcStepLo--\nDcStepHi \\
cooperativity exponent & \(n\) & -- & \nDcNexpLo--\nDcNexpHi \\
Henry-site fraction & \(f_H\) & -- & \nDcFhLo--\nDcFhHi \\
particle diameter & \(d_p\) & mm & \nDcDpLo--\nDcDpHi \\
particle density & \(\rho_p\) & kg\,m\(^{-3}\) & \nDcRhopLo--\nDcRhopHi \\
bed porosity & \(\varepsilon_t\) & -- & \nDcEpsLo--\nDcEpsHi \\
\midrule
\multicolumn{4}{l}{\emph{operating condition}} \\
feed relative humidity & \(\mathrm{RH}_{\mathrm{feed}}\) & -- & \nDcRHLo--\nDcRHHi \\
superficial velocity & \(v\) & m\,s\(^{-1}\) & \nDcVLo--\nDcVHi \\
feed temperature & \(T_{\mathrm{in}}\) & K & \nDcTinLo--\nDcTinHi \\
\bottomrule
\end{tabular}
\end{table}

\section{The surrogates of each rung}

\begin{sidewaystable}
\centering
\caption{Surrogates by rung. Absolute errors in different rows are not comparable.
POD, proper orthogonal decomposition; NTK, neural tangent kernel; L-BFGS, limited-memory
Broyden--Fletcher--Goldfarb--Shanno; DeepOKAN, a DeepONet with Kolmogorov--Arnold layers;
\(p\), basis size.}
\label{s2:tab:rungs}
\footnotesize
\renewcommand{\arraystretch}{1.0}\setlength{\tabcolsep}{4pt}
\begin{tabular}{@{}>{\raggedright\arraybackslash}p{0.08\textheight}>{\raggedright\arraybackslash}p{0.11\textheight}>{\raggedright\arraybackslash}p{0.2\textheight}>{\raggedright\arraybackslash}p{0.31\textheight}>{\raggedright\arraybackslash}p{0.19\textheight}@{}}
\toprule
rung & data, design & input \(\to\) output; encoding & architecture & training \\
\midrule
L1 & larger, five-fold &
parameters \(\to\) POD coefficients per channel &
ridge; random forest; gradient boosting; perceptron with standardised targets &
scikit-learn; perceptron Adam with early stopping \\
L1 curve & larger, five-fold & parameters \(\to\) \nLCpField\ POD coefficients of \(c\); training materials subsampled & one gradient-boosted regressor per mode (not the L1 perceptron) & three seeds per size \\
L2 & larger, five-fold & as L1 &
perceptron width and depth, boosting depth, forest leaf size, swept in both directions & as L1 \\
L3 & legacy &
parameters \(\to\) POD coefficients; full stored field &
perceptron and two Kolmogorov--Arnold \cite{liu2025kan} families (Gaussian radial basis, Chebyshev), depth \nMthThreeDepth, matched to \nMthThreeBudgets\ parameters; widths \nMthThreeWidths &
Adam, cosine schedule, \nMthThreeSteps\ steps, per-family learning-rate grid \\
L4 & larger; time axis and single split &
(parameters, \(z\), \(t\)) \(\to\) \((c,q,T)\) pointwise, bounded output head &
one perceptron, \nMthFourDepth\ layers of \nMthFourWidth\ (\nMthFourParams\ parameters); data-only twin and physics-weighted arms (fixed, gradient-norm, NTK, self-adaptive); residual of equations \eqref{eq:mass}--\eqref{eq:energy} with the Danckwerts inlet conditions, corrected (B72) &
Adam, cosine schedule, \nMthFourSteps\ steps of \nMthFourSup\ data, \nMthFourColloc\ collocation and \nMthFourBc\ boundary points; then L-BFGS (polish) \\
L5 & legacy &
parameters \(\to\) \(c\) field, subsampled from the stored field &
DeepONet and DeepOKAN over basis size \(p\); Fourier neural operator over mode count; all matched to \nMthFiveBudget\ parameters &
Adam, cosine schedule, \nMthFiveSteps\ steps, per-arm learning-rate grid \\
L5 basis size & larger, five-fold & as L5 & DeepONet over basis size \(p\), matched parameters &
as L5, step budget tripled \\
L5 basis vs.\ map & legacy; larger, five-fold &
parameters \(\to\) POD coefficients of \(c\) over basis size &
one gradient-boosted regressor per mode & one seed (legacy), three seeds (larger) \\
L6 & larger, five-fold &
\nMthSixNdesc\ observable descriptors (isotherm at three temperatures, bed and operating values), \(z\), \(t\) \(\to\) \((c,q,T)\) &
joint perceptron versus separate equilibrium, kinetic and transport heads, with \(q\) from equation~\eqref{eq:ldf}; both \nMthSixBudget\ parameters, depth \nMthSixDepth &
Adam, cosine schedule, \nMthSixSteps\ steps \\
L7 & larger, five-fold, paired with L1 &
parameters \(\to\) outlet curve &
equilibrium shock, Klinkenberg, constant pattern; nothing fitted & none \\
Warp & legacy &
parameters \(\to\) \(c\), in a fixed frame or the two-front frame of equation~\eqref{s2:eq:warp} &
POD plus one gradient-boosted regressor per mode; front trajectories from a POD of the landmark curves &
three seeds; oracle arm given the true fronts \\
\bottomrule
\end{tabular}
\end{sidewaystable}

\section{The coordinate change, and where the obstruction moved}

\label{s2:sec:warp}

\begin{figure}[!tbp]
\centering
\includegraphics[width=\textwidth]{../figures/Fig6_warp.pdf}
\caption{The two-front frame (legacy design). \textbf{A}~Modes needed for 99.9\,\% of
training variance in each frame: aligning one front raises the count, aligning both
lowers it. Tick labels give the level of \(c/c_{\mathrm{in}}\) at which each aligned
front is located (one level, or a pair); the hatched pair was excluded by its
interpolation guard. \textbf{B}~Held-out error: with predicted fronts the two-front frame is not
distinguishable from the fixed frame (the interval runs from
\nWarpPredImproveMax\,\% better to \nWarpPredDegradeMax\,\% worse, \nWarpNmat\
materials); handed the true fronts (hollow) it is better, an upper bound rather than a
result.}
\label{s2:fig:warp}
\end{figure}

\subsection{Scope and prior art}

Aligning multiple transports is not new, and none of it is claimed here. The
two-transport case is the founding example of the shifted POD \cite{reiss2018shifted},
multi-front variants follow
\cite{nair2019transported,mendible2020dimensionality,krah2025robust,krah2023front,zorawski2024neural},
a neural-network shifted decomposition predates the last of these
\cite{papapicco2022nnspod}, and breakthrough curves have been decomposed with a single
characteristic time \cite{lee2026scaleup}. Our two-landmark linear warp is piecewise
landmark registration from functional data analysis \cite{kneip1992statistical} (a
landmark-free alternative is in \cite{ramsay1998curve}); Zucatti and Zahr
\cite{zucatti2026model} apply the same primitive to shock features in reduced-order
models of convection-dominated problems, and registration-based reduction is an
established family \cite{taddei2020registration,taddei2026erratum}. What is ours is the
measurement.

\subsection{One shift makes it worse; two make it much better}

The standard remedy for a transport-dominated \(n\)-width is to factor the front
position out before projecting \cite{reiss2018shifted}. We warped each field onto its own local arrival time
and swept the alignment level. \emph{It makes the \(n\)-width worse}: aligning on
the half-height crossing takes the modes needed for 99.9\,\% of training variance
from \nNwidthOriginal\ to \nNwidthSingle. It does not help accuracy even given
the exact front trajectory for free: handed the oracle warp, the co-moving frame
still loses to the fixed frame. The reason was measured before it was invoked: the separation between the two waves spans a factor of \nGapRange, and one shift cannot align
two waves whose separation varies that much.

Using \emph{two} shifts, one per wave, and normalising the interval between them (Figure~\ref{s2:fig:warp}),
\begin{equation}
\sigma(z,t) = \frac{t - t_{\mathrm{lo}}(z)}{t_{\mathrm{hi}}(z) - t_{\mathrm{lo}}(z)},
\label{s2:eq:warp}
\end{equation}
collapses the \(n\)-width where one shift inflated it: \nNwidthOriginal\ \(\to\)
\nNwidthTwoWaveA\ modes, or \nNwidthTwoWaveB\ for the wider landmark pair. A third
pair was excluded because a narrow window divides by a small span and its
interpolation floor reached \nComovingFloorPct\,\% of the fixed-frame error it is judged
against, above the guard's threshold of \(10\,\%\).

\subsection{The verdict, what it can resolve, and the relocation}

Table~\ref{s2:tab:warp} gives the result: three seeds for every stochastic component, both
arms computed in the same run, and a paired cluster bootstrap by material.

\begin{table}[t]
\centering
\caption{The two-front frame against the fixed frame (legacy design). Differences are
the fixed frame's nRMSE minus the arm's, so positive means the arm is better; CI is the
cluster-bootstrap confidence interval. The oracle arm is handed the true front
trajectories and is \emph{not} comparable with any arm that is not; it sizes the
remaining headroom and is not a result.}
\label{s2:tab:warp}
\resizebox{\textwidth}{!}{%
\begin{tabular}{lrl}
\toprule
arm & held-out nRMSE & verdict vs.\ the fixed frame \\
\midrule
fixed frame & \nWarpFixed & reference \\
two-front, predicted fronts & \nWarpPredicted & \nWarpPredDiff, CI [\nWarpPredLo, \nWarpPredHi]: no difference \\
two-front, oracle fronts & \nWarpOracle & \nWarpOracleDiff, CI [\nWarpOracleLo, \nWarpOracleHi]: significant \\
\bottomrule
\end{tabular}}
\end{table}

The predicted-front arm does not beat the fixed frame. Given the true front
trajectories, the same POD-plus-regressor arm is \nWarpHeadroomFixed{}\(\times\) better
than the fixed frame, which is what perfect front location would buy, and
\nWarpHeadroom{}\(\times\) better than with predicted fronts. On \nWarpNmat\ held-out
materials the predicted-front interval admits anything from a \nWarpPredImproveMax\,\%
improvement to a \nWarpPredDegradeMax\,\% degradation; it rules out the oracle's gain,
not a modest one. With that scope, this is a relocation of the obstruction. Through L1--L5 the obstruction was the
parameter-to-coefficient map, at the number of modes a given training set supports. It
is now ``two smooth monotone one-dimensional curves are not located accurately
enough'', a better-posed problem, with a target a front-locating model can be held
to. The oracle's gain is not unique to this route: the Fourier operator of L5 reaches a
comparable error on the same held-out materials without any oracle.

\paragraph{Monotonicity is not the lever} The predicted fronts violate the
constraints the true ones satisfy in \nWarpNonmonoLo\,\% of samples, but projecting
them onto the constraint set makes the reconstruction slightly worse
(Supplementary S2, Section S8); the remaining headroom is the front predictor's raw
accuracy.

\paragraph{Relation to constant-pattern behaviour} Equation~\eqref{s2:eq:warp} is not
constant-pattern scaling \cite{knox2016limitations}: when the two fronts move at
different rates and their separation varies \nGapRange-fold, no single scale factor
aligns both, which is why the second landmark does the work.

%% ===========================================================================

\section{The accuracy--cost frontier}

\begin{table}[t]
\centering
\caption{The accuracy--cost frontier, in solve-equivalents. Each row is a point on
the learning curve of Section~\ref{sec:ladder}, priced; its error is that of the curve's
per-mode gradient-boosting learner, and the fit and query times quoted in the text are
those of the L1 perceptron. The last column is a \emph{lower} bound on the number of
screening queries at which the surrogate becomes cheaper than the reference; it needs
only the data cost, so it holds for either learner and whatever a query costs.}
\label{s2:tab:cost}
\resizebox{\textwidth}{!}{%
\begin{tabular}{rrrr}
\toprule
training materials & training simulations & held-out nRMSE & break-even queries \(\geq\) \\
\midrule
12  & \nCostSmallSims & \nLCsmallest & \nCostSmallBreakEven \\
24  & \nCostSimsB     & \nCostErrB   & \nCostBeB \\
48  & \nCostSimsC     & \nCostErrC   & \nCostBeC \\
96  & \nCostSimsD     & \nCostErrD   & \nCostBeD \\
192 & \nCostLargeSims & \nLClargest  & \nCostBreakEven \\
\bottomrule
\end{tabular}}
\end{table}

Second, the learning curve is a price list: going from 12 to 192 training materials
multiplies the data cost by \nCostDataRatio{}\(\times\) and divides the error by
\nLCratio{}\(\times\). The rows are priced for the curve's per-mode regressor; the L1
perceptron is already more accurate, and L2 found a configuration
\(\nLtwoVsLonePct\,\%\) better still, so a better learner shifts the whole frontier
down. The exchange rate between data and error is the transferable part.

Third, the eliminated rungs are interventions that add cost without buying accuracy: a
Kolmogorov--Arnold map costs fit time and nothing else, and the closed forms are
cheaper but worse even when fitted. Of the levers that remain, only training materials
cost simulations; a better configuration, a small residual weight
(\nLfourbMatGainPct\,\%) or a different operator cost fit time or collocation points.
Every statement here is cost \emph{at a stated accuracy}: the reference answers exactly,
the surrogate with the held-out error in the third column.

%% ===========================================================================

\section{Corrections and retractions}

\label{s2:sec:corrections}

\nLedgerA\ claims made in this project were withdrawn and \nLedgerB\ further defects
were caught and corrected; Supplementary S1 lists every entry, and the text cites them
by code (A for a withdrawn claim, B for a defect). We report them because six reversed
a headline verdict, and a reader who does not know which numbers moved cannot evaluate
the ones that did not; disclosed self-correction is also what lets a reader judge the
rest \cite{rohrer2021loss}.

\paragraph{Six that reversed a headline} Two reversed the physics-as-a-loss rung: its
legacy claim that the residual hurts did not survive a sweep of weightings (B69), and
every one of its verdicts was then withdrawn when the residual proved to be the wrong
equation (B72, A28), until the corrected re-run reported the small benefit of
rung L4 (Section~\ref{sec:ladder}). A17 withdrew the explanatory role of the Kolmogorov
\(n\)-width; A22 turned a null described as decisive into one that excludes only very
large effects (A25 later corrected its power simulation in the opposite direction);
A29 retracted the flat-in-basis verdict of L5 on its \nNclustFolds-material
re-measurement; and A30 revised the cost case once the coarse-solver baseline proved
more accurate than the perceptron surrogates. B75, which removed ``significantly'' from
the Fourier operator's advantage over DeepONet once the choice of both arms was paid
for, is not counted: the direction of that comparison stands.

The errors recur in five classes: a property measured at one sample size and stated as
intrinsic (A18, A21, A26, A29); normalising by a quantity that can vanish (A15, B34,
B41); comparison against a weak configuration of the competitor (A20, B14, B22, B23);
undertrained arms scored as architectures (A19, B24); and power and calibration (A13,
A22, A25).


%% ===========================================================================

